
% ============================================================================
\clearpage
\section*{v4.9.0-r85 Source-Only: Masses and Couplings in the Second Reading, the Epoch \texorpdfstring{\(\tau^*\)}{tau*}, and the Big Bang as a CPT Transition}
\addcontentsline{toc}{section}{v4.9.0-r85 Source-Only: Masses and Couplings in the Second Reading, the Epoch tau*, and the Big Bang as a CPT Transition}
\label{sec:v490-r85-masses-cpt}
% ============================================================================

\paragraph{Purpose.}
Layer v4.9.0-r84 fixed the tachyon part as a second reading of the same field, formed by the complete complex boost with the
holomorphic adjoint (\S\ref{sec:v490-r84-second-reading}).  On 5 November 2026 the author asked to check the mass terms, the weighting
of gravity, and the two centres of the time reflection further.  This layer records five pre-registered tests (H82--H86, all not blind),
two explorations, one author decision (the domain G2) and one working assumption.  The notes are
\texttt{notes/\allowbreak research/\allowbreak v4\_10\_0\_r85/}.  Earlier layers are not rewritten.

\noindent\textbf{Source-only status.}
This v4.9.0-r85 layer is source-only.  It does not build a PDF, does not assign a new DOI,
and does not authorize a public release.  The v4.9.0 Version DOI
\href{https://doi.org/10.5281/zenodo.23534954}{10.5281/zenodo.23534954} and the Concept DOI
\href{https://doi.org/10.5281/zenodo.19210728}{10.5281/zenodo.19210728} are unchanged by it.

\subsection*{v4.9.0-r85 results}
\label{subsec:v490-r85-results}

\paragraph{Theoremlet v4.9.0-r85.1 (H82: masses and couplings).}
Pre-registered (commit \texttt{c3f956c}, not blind).  With \(S=-i\sigma_z\), \(S^{\mathsf T}\varepsilon S=\varepsilon\): Majorana and Dirac mass terms are
invariant under the complete boost and enter with \(P=a+b\); terms with derivatives enter with \(K=a-b\).  After canonical normalisation, at
matter weights (\(K=1/\gamma\), \(P=1\)): \(\alpha\), fermion and Majorana masses and \(G_N\) scale with \(\gamma\); \(m_H\) and \(\ell_P\) with
\(\sqrt\gamma\); \(v\) with \(\gamma^{-1/2}\); \(m_W\) and \(\theta_W\) are unchanged (a, confirmed).  Compared with the weights of v4.9.0-r19, \(\alpha\),
\(\alpha_s\) and \(\theta_W\) agree, \(G_N\) (\(1\) against \(k=2\)) and \(\ell_P\) (\(\tfrac12\) against \(\tfrac32\)) do not (b, refuted as expected).  The
different growth follows a fixed pattern given by the number of derivative terms; all factors can be absorbed into redefined parameters.

\paragraph{Theoremlet v4.9.0-r85.2 (H83: three rules for gauge fields and gravity).}
Pre-registered (commit \texttt{36ec8e7}, not blind).  R1 as in r85.1; R2 computed from \(g\mapsto-g\) with fixed fields (reading B for a scalar
is exactly this), \(R[-g]=-R[g]\), \(\sqrt{|g|}\) unchanged; R3 a single action on \(c\,g\).  R2 gives Einstein--Cartan \(-1\), Yang--Mills \(+1\)
(two inverse metrics), Weyl kinetic term \(-i\) (one inverse vierbein).  No rule gives \(k(\alpha)=1\) and \(k(G_N)=2\) together; \(G_N\) scales
with \(\gamma^1\) under all three.  The match for \(\alpha\) in r85.1 depends on R1.

\paragraph{Exploration v4.9.0-r85.3 (the phase of the spinors).}
Under R2 the fermion action is \((a-ib)S_{\rm Weyl}\): in the real reading its tachyon image is invisible, in the complex reading K
\(\operatorname{Im}S=-b\,S_{\rm Weyl}\) forces the energy--momentum tensor and the spin density of the fermions to vanish.  Quantum mechanically
\(a-ib\) is a partial Wick rotation by \(\arctan(b/a)\) (\(16.8^\circ\) at matter weights, \(45^\circ\) at photon weights).  \(g\mapsto-g\) is not
unique for spinors, since \(Cl(1,3)\not\cong Cl(3,1)\); F3 allows the phase \(\mp i\) for a weight \(\tfrac12\).  \emph{Working assumption}
(author, ``carry on'' on the proposal): the Weyl kinetic term has F3 weight \(1\), as in r81 and r84.3; not an explicit decision.

\paragraph{Theoremlet v4.9.0-r85.4 (H84: quantum modes at \(\tau^*\)).}
Pre-registered (commit \texttt{d2a19f2}, not blind).  For the offset \(\delta\neq0\) of r84.4 the Wronskian \(K(\chi\bar\chi'-\bar\chi\chi')=i\) stays constant
(\(4\times10^{-14}\)); since \(K(\tau^*)=0\), every mode with this normalisation has \(\chi'\propto1/|\tau-\tau^*|\), and its energy density
\(E\,|\tau-\tau^*|\to-0.1496\) from below and \(+0.1494\) from above.  The mode selection of r84.4(c) is classical only; for a quantum field
the offset is singular at \(\tau^*\), which favours the reflection about the big bang.

\paragraph{Author decision (5 November 2026): the domain G2.}
The reflection maps the domain as well: the image of \(\tau>0\) is \(\tau<0\).  r84.4(a) (H80) took both readings on the whole time axis
(domain G1).  The author chose \textbf{G2}: the matter reading lives after the big bang, the tachyon reading on its image before it, joined
only at the big bang.  This is structurally the two-sheet picture of the CPT-symmetric universe of Boyle, Finn and Turok
(arXiv:1803.08928).  r84.4(a) stays valid for G1.

\paragraph{Exploration v4.9.0-r85.5 (the sharp transition under G2).}
For conformal fields the weight jumps from \(-b\) to \(a\) at \(\tau=0\).  With continuity of the field and of weight times derivative, the
positive-norm mode before the big bang becomes \(C e^{-i\omega\tau}+De^{+i\omega\tau}\), \(C=(a+b)/(2a)\), \(D=(a-b)/(2a)\); normalised
\(\alpha=(1+x)/(2\sqrt x)\), \(\beta=(1-x)/(2\sqrt x)\) with \(x=\bsym^2\), \(\beta/\alpha=1/\gamma=0.536\), \(n=(1-x)^2/(4x)=0.404\) per mode for every \(k\)
(sympy).  At photon weights \(n=0\).

\paragraph{Theoremlet v4.9.0-r85.6 (H85, H86: the smooth transition).}
A real smooth weight from \(-b\) to \(a\) has a zero and is singular by r85.4; the weight is therefore taken around zero in the complex
plane, \(w=b^{1-s}a^se^{\pm i\pi(1-s)}\), a partial Wick rotation, with conserved holomorphic Wronskian.
\begin{enumerate}[leftmargin=2em]
\item[(a)] H85 (commit \texttt{4ac77a0}, partly not blind; \(\Delta=1\), tanh profile): both quantitative criteria missed.  \(D\) approaches the
sharp value only linearly in \(k\Delta\) (\(3.5\%\) at \(0.01\)), and the pre-registration had exchanged the roles of \(C\) and \(D\).
\item[(b)] Noticed afterwards: adiabatically the positive-norm mode of the tachyon reading continues as \(e^{+ik\tau}\), which after the big bang
is the mode of negative frequency; the admixture \(C\) falls off exponentially; the two senses of rotation distribute \(D\) unequally.
\item[(c)] H86 (commit \texttt{851297a}, not blind, motivated by (b); new settings \(\Delta=2\), tanh and error-function profiles): confirmed.
For \(k\Delta\ge6\), \(|C|<10^{-2}\) and \(\sqrt{|D_+||D_-|}=\sqrt{b/a}\) to \(10^{-13}\); \(\ln|C|\) is linear on \(k\Delta\in[2,8]\) (\(R^2=0.99987\)) with slope
\(-2.67\), within the \(15\%\) criterion around \(\pi\) only narrowly and not tending to \(\pi\) in the range; at \(k\Delta=0.01\), \(C\) meets the sharp
value to \(2\times10^{-4}\); at \(k\Delta=0.02\), \(|D_-|/|D_+|=1.13\) (tanh) and \(1.10\) (erf) with \(\sqrt{|D_+||D_-|}/D_0=0.998\) and \(0.999\).
\end{enumerate}

\paragraph{Proof and measurement sources.}
Scripts \texttt{k197}--\texttt{k201} and the notes in \texttt{notes/\allowbreak research/\allowbreak v4\_10\_0\_r85/}; pre-registrations and results
H82--H86 in \texttt{notes/\allowbreak research/\allowbreak predictions/}.  \(\square\)

\paragraph{Consequence.}
Mass terms carry over to the second reading; the pattern of factors matches v4.9.0-r19 for the gauge couplings but not for gravity.
For a quantum field the reflection about the big bang is preferred over an offset.  Under G2 with a complex transition the big bang maps
every mode of the tachyon reading before it to a mode of negative frequency after it (a CPT assignment), with an admixture fixed by
\(\bsym\) for long waves and cut off exponentially above \(1/\Delta\); \(\Delta\) is a free scale.  The two senses of the Wick rotation give an
asymmetry with fixed product; its physical meaning is open.  OP5 is not closed.

\subsection*{v4.9.0-r85 status in the three frames}
\label{subsec:v490-r85-frames}

\begin{itemize}[leftmargin=2em]
\item Under G2 the matter frame M holds after the big bang and the tachyon frame T before it; the photon frame P is the equal-weight
mixture, where the transition of r85.5 produces no particles (\(n=0\)).  This is a statement at P in the sense of F2.
\item The factors of r85.1 and the numbers of r85.4--r85.6 use the weights of frame M.  The comparison with v4.9.0-r19 relates a reduction
(M and T to one field) to a conversion (M to P); these are different operations.
\end{itemize}

\subsection*{v4.9.0-r85 register entries}
\label{subsec:v490-r85-register}

\begin{center}
\small
\begin{tabular}{@{}>{\raggedright\arraybackslash}p{0.34\textwidth}>{\raggedright\arraybackslash}p{0.58\textwidth}@{}}
\toprule
statement & status after v4.9.0-r85\\
\midrule
masses in the second reading & carry over; factors \(\gamma\), \(\sqrt\gamma\), \(1\) by derivative count (H82)\\
weights of v4.9.0-r19 from the reduction & gauge couplings yes, gravity no, under three rules (H82, H83)\\
F3 weight of the Weyl kinetic term & \(1\) (working assumption); \(\tfrac12\) would hide or forbid fermions\\
offset \(\delta\neq0\), quantum & singular at \(\tau^*\) (H84)\\
domain of the readings & G2: T before, M after the big bang (author decision)\\
big bang under G2 & CPT assignment plus admixture fixed by \(\bsym\), cut off above \(1/\Delta\) (H86)\\
\bottomrule
\end{tabular}
\end{center}

\subsection*{v4.9.0-r85 ledger and guard update}
\label{subsec:v490-r85-ledger-guard}

\[
R_{\rm chk}^{v490\text{-}r85}=(1,1,0)
\]
for (masses and the r19 comparison decided; the big-bang transition decided for G2; scale \(\Delta\) not fixed).

The layer spends no guard:
\[
R_{\rm glob}^{(v490\text{-}r85,+)}=(0,1,1,1,1),
\qquad
\Gamma_{\rm guard}^{v490\text{-}r85}=(1,1,1,1,1),
\qquad
\beta_4=1/256\ \text{unspent}.
\]
It does not imply exact finite \(C5\) zero closure, exact all-readout terminality,
\(C6/T2\), a global photon-clock/RH theorem, or global all-readout closure.
